% Transcription — fonds Alexandre Grothendieck, shelfmark 16, batch 3 (pages 41-55).
% Established from the facsimile archives/batches/16/batch-03.pdf (a mirror of
% https://grothendieck.umontpellier.fr/16.pdf), per the skill
% transcribe-grothendieck. The batch file carries no cover sheet: its PDF page 1
% is archivists' page 41 and its PDF page 15 is archivists' page 55, checked
% against the pencilled numbers at the bottom-left (41, 43, 44, 45, 46, 47, 48,
% 49, 50, 52 are legible on their leaves).
%
% Pass: Opus 5 (claude-opus-5), 2026-09-19 — first pass, unchecked against the pages by a human.
%
% What the batch is. The folder's last fifteen leaves, and the only batch of the
% three that carries correspondence. Eleven pages are transcribed: 41, 50, 52
% and 54 in his hand, and 43-49, two carbon copies of letters to Coates, typed
% in English with Greek letters, arrows and corrections inked by hand. The other
% four are skipped: 42 is a galley of EGA IV (§ 18.11, unramified homomorphisms),
% unrelated to this folder's subject although it carries a few corrections of
% his; 51, 53 and 55 are a Bourbaki « Tribu » n° 65 on étale algebras and the
% different, and carry nothing of his. Page 54 stops a third of the way down and
% the rest of the leaf is blank; the folder ends there.
%
% The two letters are dated in his own hand at the head, « 4.1.67 » on page 43
% and « 6.1.1966 » on page 47 — that is what the leaves carry, and the
% transcription records it without trying to reconcile the two years.
%
% The mathematics is one subject seen from three sides: the standard
% conjectures A, B, C, D and what makes them independent of the polarization and
% of the cohomology theory.
%
%   41   a working sheet, in his hand, on the equivalence « moyennant C(Y),
%        C(X) équivaut à : Λ_X φ_* est algébrique ». The class of Λ_X φ_* lives
%        in Σ H^{2(n-1)-i}(Y) ⊗ H^i(X) = H^{m-1}(X×Y); acting on it by
%        L_{X×Y} = L_X ⊗ 1 + 1 ⊗ L_Y gives
%        Λ_X φ_* L_Y + L_X Λ_X φ_* = (Λ_X L_X + L_X Λ_X) φ_*
%                                  = (φ_* Λ_Y φ^* + id_X) φ_*.
%        That identity is exactly the one the letter of page 45 quotes, so this
%        leaf is the scratch behind it.
%   43-46 letter of 4.1.67: where to put the index conjecture in Coates's notes;
%        in characteristic zero A_χ to D_χ are independent of χ and A and C of
%        the polarization, so they can be written A(X) to D(X); the same is not
%        known in characteristic p. Then a Proposition in French, (i) to (iv bis)
%        plus a Corollaire (v), (vi), giving conditions equivalent to C_χ(X) in
%        terms of hyperplane sections, and the formula
%        (Λ_X φ_*)L_Y + L_X(Λ_X φ_*) = (φ_* Λ_Y φ^* + id_X) φ_*. A list of what
%        is known (dim X ≤ 2 in any characteristic; π_0, π_{2n}, π_1, π_{2n-1}
%        algebraic; Hodge 1938 on surfaces; A(X) for dim X ≤ 4 in car. 0), and
%        the remark that everything already fails for 1-cycles on threefolds.
%   47-49 letter of 6.1.1966: a Proposition making C_χ(X) equivalent to a)
%        D_χ(X) plus an algebraic isomorphism H^{2n-i}(X) ≅ H^i(X) for i < n, and
%        to b) rationality of the characteristic polynomials of algebraic
%        endomorphisms of H^i(X) plus the same isomorphisms. The proof runs
%        through Hamilton-Cayley: w = uv is an algebraic automorphism, w^{-1} is
%        a rational combination of the w^i, hence algebraic, hence so is v^{-1}.
%        A Corollary on u : H^i(X) → H^{i+2j}(X'), the « complete yoga » that the
%        category of motives is semi-simple, and the weak variants A'_χ, A''_χ
%        which in fact imply the full Lefschetz relations for algebraic cycles.
%   50, 52, 54 the « old scribbles » the letter of page 46 promises to look at
%        again: a Lemme on a ring A of endomorphisms of H*(X), stable under
%        homogeneous components and transposition, containing ζ, and such that
%        each ζ^{n-i} has a left inverse u_i in A — then A contains L and the
%        projections on the elementary factors ζ^j P^i. The proof is by
%        induction on k = i + 2j, through T_k = id − Σ_{i<k} p_i − Σ_{j>2n-k} p_j.
%
% Notation fixed for the batch: L_X and Λ_X for the Lefschetz operators, drawn
% as a bracket and a wedge; φ for the inclusion of the hyperplane section Y in
% X, with φ^* and φ_*; p_i for the projection on H^i; P^i for the primitive
% part; ζ for the operator the Lemme is given (an L without its name); χ for the
% index of the cohomology theory and ξ for the polarization, both inked by hand
% into the typescript and sometimes left out by him — where a subscript is
% missing on the leaf, the transcription leaves it missing. The typescripts'
% misspellings (« independant », « arbitray », « correspondance », « whos »,
% « week » for weak, « Condititions », « décorissante ») are his and stay.
\documentclass[11pt,a4paper]{article}
% Le préambule commun à toutes les transcriptions du fonds Grothendieck.
%
% Il tient en une idée : **l'incertitude se compose, elle ne se résout pas.**
% Une transcription qui lisse un mot douteux a détruit la seule chose pour
% laquelle on la faisait — un lecteur doit pouvoir savoir, à chaque endroit,
% ce qui était sur le feuillet et ce qui a été ajouté. Les six macros qui
% suivent sont donc l'appareil critique, et non de la décoration.
%
% Les mêmes commandes sont reconnues par `scripts/render.mjs`, qui produit la
% vue de lecture du site. Une macro ajoutée ici doit l'être aussi là-bas,
% faute de quoi le rendu échouera bruyamment — ce qui est voulu : un rendu
% silencieusement faux est pire qu'un rendu absent.

\ProvidesPackage{grothendieck}[2026/08/08 Transcriptions du fonds Grothendieck]

\usepackage{amsmath,amssymb,amsthm}
\usepackage{mathrsfs}
\usepackage{stmaryrd}
% Les diagrammes commutatifs sont partout dans ce fonds : c'est la langue
% même de la géométrie algébrique de Grothendieck, et une transcription qui
% les rendrait en prose ne transcrirait rien.
\usepackage{tikz-cd}
% Français par défaut — c'est la langue des feuillets — mais l'anglais est
% chargé aussi : la traduction et le résumé sont des documents anglais, et la
% typographie française (espaces avant les deux-points) y serait une faute.
% Ces documents basculent par \selectlanguage{english} en tête de corps.
\usepackage[english,french]{babel}
\usepackage{fontspec}
\usepackage{geometry}
\geometry{a4paper,margin=25mm,marginparwidth=28mm}
\usepackage{marginnote}
\usepackage{xcolor}
% ulem, not soul. `soul`'s \st and \ul are built on a character-by-character
% scan that XeLaTeX's font machinery breaks: the document fails with an
% undefined control sequence rather than degrading. ulem does the same two jobs
% and is Unicode-engine safe. `normalem` stops it hijacking \emph, which the
% transcriptions use for Grothendieck's own emphasis.
\usepackage[normalem]{ulem}
\usepackage{hyperref}
\hypersetup{colorlinks=true,linkcolor=black,urlcolor=black,pdfborder={0 0 0}}

% --- Métadonnées du lot -----------------------------------------------------
% Reprises telles quelles par le script de rendu, qui les affiche en tête de
% la vue de lecture. Elles fixent ce que le fichier prétend couvrir : sans
% elles, une transcription flotte sans point d'attache dans un fonds de seize
% mille feuillets.

\newcommand{\folder}[1]{\gdef\grothFolder{#1}}
\newcommand{\batch}[1]{\gdef\grothBatch{#1}}
\newcommand{\leaves}[2]{\gdef\grothFirst{#1}\gdef\grothLast{#2}}
\newcommand{\foldertitle}[1]{\gdef\grothFoldertitle{#1}}
\gdef\grothFolder{?}\gdef\grothBatch{?}\gdef\grothFirst{?}\gdef\grothLast{?}\gdef\grothFoldertitle{}

% --- Le résumé --------------------------------------------------------------
%
% Il ouvre la lecture modernisée, sous le titre « Résumé », et il est composé
% en petit. Ce n'est pas un ornement : c'est le seul passage du document qui
% ne soit pas des mathématiques, et un lecteur doit pouvoir le reconnaître
% avant de l'avoir lu — puis le sauter s'il connaît déjà le sujet, ou s'y
% arrêter s'il ne le connaît pas. Le filet qui le suit marque la reprise du
% corps.

\newenvironment{resume}
  {\small\setlength{\parskip}{0.45em}}
  {\par\medskip\hrule\medskip}

% --- Mots-clés ---------------------------------------------------------------
%
% En anglais, en clôture du résumé de la lecture modernisée : le vocabulaire
% moderne sous lequel chercher ce que les feuillets construisent. Le manifeste
% du site (`npm run manifest`) les extrait pour étiqueter la cote entière —
% cette ligne est la source unique des tags, il n'y a pas de fichier à côté.
\newcommand{\keywords}[1]{\par\smallskip\noindent
  {\footnotesize\textsc{Keywords} --- \textit{#1}}\par}

% --- L'appareil critique ----------------------------------------------------

% Le numéro de feuillet, en marge. C'est l'ancre qui permet de régler un
% désaccord sur une formule en regardant *un* feuillet plutôt qu'une chemise
% de six cents. Le site s'en sert aussi pour tourner les pages du fac-similé
% quand on fait défiler la transcription.
\newcommand{\leaf}[1]{%
  \marginnote{\footnotesize\color{gray!70}\textsf{#1}}%
  \label{leaf:#1}\ignorespaces}

% Illisible : on ne devine pas. Un mot inventé qui se lit comme les autres est
% le pire résultat possible.
\newcommand{\ill}{\textcolor{red!60!black}{[\,\ldots\,]}}

% Lecture douteuse : le mot est donné, et signalé comme incertain.
\newcommand{\uncertain}[1]{\uline{#1}}

% Ajout de l'éditeur — une lettre manquante, un mot sous-entendu.
\newcommand{\add}[1]{\textcolor{blue!50!black}{[#1]}}

% Biffé par Grothendieck lui-même. Ce qu'il a rayé fait partie du document :
% une rature dit souvent où la pensée a changé de direction.
\newcommand{\struck}[1]{\sout{#1}}

% Note du transcripteur, sur l'état matériel du feuillet.
\newcommand{\note}[1]{\par\smallskip{\small\color{gray!60!black}\textit{[#1]}}\par\smallskip}

% Note marginale de Grothendieck, distincte du corps du texte.
\newcommand{\marginal}[1]{\par\smallskip\noindent\hspace{1em}%
  {\small\color{gray!60!black}#1}\par\smallskip}

% --- Titre ------------------------------------------------------------------

\AtBeginDocument{%
  \begin{center}
    {\large\bfseries Fonds Alexandre Grothendieck --- cote n\textsuperscript{o}~\grothFolder}\\[2pt]
    {\itshape\grothFoldertitle}\\[4pt]
    {\small feuillets \grothFirst--\grothLast\ (lot \grothBatch)}\\[2pt]
    {\footnotesize\color{gray!60!black}Transcription établie d'après le fac-similé numérisé
     par l'Université de Montpellier.\\
     Les lectures incertaines sont soulignées, les illisibles notées [\,\ldots\,],
     les ajouts entre crochets.}
  \end{center}
  \vspace{1em}}

\endinput


\folder{16}
\batch{3}
\pages{41}{55}
\dating{1966-1967}
\watermark{Édition de démonstration}
\foldertitle{Formalisme algébrique des correspondances et algèbres de Poincaré-Hodge (cf conjectures standard) : notes manuscrites (s.d.), tapuscrit annoté (s.d.), lettres (1966-1967)}

\begin{document}

\section*{Note manuscrite : $C(X)$ et l'algébricité de $\Lambda_{X}\varphi_{*}$ (page 41)}

\note{La page s'ouvre au milieu d'un raisonnement : elle n'a pas de titre et
reprend la discussion de $C(X)$ et de $C(Y)$ menée au « Formulaire » des
pages 32 à 39, que le bas de la page 39 laisse ouverte sur un « Mais ». $X$ est de
dimension $n$, $Y \subset X$ une section hyperplane, $\varphi : Y \to X$
l'inclusion ; $m = \dim X \times Y = 2n-1$.}

\page{41}

Moyennant $C(Y)$, $C(X)$ équivaut à la condition :

$\forall i \leq n-1$, l'hom.\ $\psi^{i}_{\xi} : H^{i}(Y) \longrightarrow
H^{i}(X)$, inverse à gauche de $\varphi^{i} : H^{i}(X) \to H^{i}(Y)$,
\struck{\ill{}} égal : $\Lambda_{X}\varphi_{*} \mid H^{i}(Y)$), est induit par
une classe algébrique.
\note{Cette condition est encadrée par un trait vertical à sa gauche. L'indice
de $\psi^{i}$ est une lettre bouclée qu'on lit $\xi$, celle-là même que le
tapuscrit des pages 43 et suivantes note pour la polarisation. Après le trait de
restriction, la parenthèse fermante n'a pas d'ouvrante.}

ou encore :

$\Lambda_{X}\varphi_{*}$ est algébrique.
\note{Également encadrée.}

\struck{ou encore} N.B.\ $\Lambda_{X}\varphi_{*}$ \struck{correspond} à une
classe de cohomologie sur $X \times Y$, qui \ill{} dans
\[
  \sum_{i} H^{2(n-1)-i}(Y) \otimes H^{i}(X) = H^{2(n-1)}(X \times Y)
  = H^{m-1}(X \times Y), \qquad m = \dim X \times Y .
\]
$L_{X} \otimes 1 + 1 \otimes L_{Y}$. Ça \struck{$\zeta_{X} \otimes 1 \,+$}
$L_{X \times Y}(\Lambda_{X}\varphi_{*})$ \ill{}
\note{Les deux mots après « qui » ne se lisent pas ; l'égalité
$\sum_{i} H^{2(n-1)-i}(Y) \otimes H^{i}(X) = \dots$ est écrite en interligne,
au-dessus de la ligne courante. $L_{X} \otimes 1 + 1 \otimes L_{Y}$ est
l'écriture de $L_{X \times Y}$.}

\[
  \Lambda_{X}\varphi_{*}L_{Y} + L_{X}\Lambda_{X}\varphi_{*}
  = (\Lambda_{X}L_{X} + L_{X}\Lambda_{X})\,\varphi_{*}
  = (\varphi_{*}\Lambda_{Y}\varphi^{*} + \mathrm{id}_{X})\,\varphi_{*}
\]
\note{Une première rédaction du second membre est biffée :
$\bigl(2\,\mathrm{id} - \sum_{n \leq i \leq 2n} p^{X}_{i}\bigr)\varphi_{*}$. Les
deux termes de gauche sont soulignés chacun d'une accolade ; sous la première il
écrit $\Lambda_{X}L_{X}\varphi_{*}$, après un essai raturé.}

(car $p^{X}\varphi_{*} = 0$)
\note{Un indice porté sous $p^{X}$ ne se lit pas ; la ligne biffée ci-dessus le
faisait courir sur $n \leq i \leq 2n$.}

Donc \ill{} algébriques. Donc \ill{}

\[
  C(X) \Longleftrightarrow \bigl[\,C(Y) \text{ et } A^{\circ}(X \times Y)\,\bigr]
\]

puis en \ill{} \ill{}

\[
  C(X) \Longrightarrow \bigl(A^{\circ}(X \times Y) \text{ et }
  A^{\circ}(Y \times Z) \text{ et } A^{\circ}(Z \times T) \dots\bigr)
\]

\struck{\ill{}} $X \supset Y \supset Z \supset T \dots$

\marginal{\ill{} discussion de \ill{} \ill{} \ill{} singulières.}

\section*{Lettre à Coates, 4 janvier 1967 (pages 43 à 46)}

\note{Carbone d'une lettre dactylographiée en anglais, sans signature ; la date
« 4.1.67 » est portée de sa main en haut à droite de la page 43. Les lettres
grecques et quelques flèches manquent à la frappe et sont encrées à la main ;
là où il ne les a pas encrées, la transcription les laisse manquer. Ses propres
mots restent en anglais ; la Proposition et son Corollaire sont en français dans
le tapuscrit. Les fautes de frappe sont les siennes et ne sont pas corrigées.}

\page{43}

Dear Coates,

I want to add a few more comments to the talk on algebraic cycles and to what I
told you on the phone.

I think the best will be to state the index conjecture right after the statement
of the main results of Hodge theory, adding that this conjecture will take its
whole significance only when coupled with « conj.A » in the next paragraph. This
will give more freedom in the next paragraph to express some extra relationships
between various conjectures, such as $A + \text{index}$ implies $B$.

In car.\ zero, state some known extra features: index theorem holds, the
\struck{conjectures} propositions $A_{\chi}$ to $D_{\chi}$ are independant of
$\chi$ (because of the existence of Betti cohomology, so that these
\struck{conjectures} propositions are equivalent to the corresponding one's for
rational cohomology), \struck{and} $A$ and $C$ are independant of the choosen
polarization \struck{$x$} (for $A$ because it is equivalent with $B$, for $C$
because it can be expressed in terms of $A$,
$C(X) = (A(X \times X) + A(Y \times Y) + \dots)$
\note{« propositions » est ajouté de sa main au-dessus de « conjectures », aux
deux endroits. La parenthèse ouverte devant $A(X \times X)$ n'est pas refermée.}

Thus the \struck{conjectures} conditions without ambiguity can be called $A(X)$
to $D(X)$, without \struck{index} subscript $\chi$ and without indication of
polarization.
\note{« conditions » et « subscript » sont frappés à la machine au-dessus des
mots qu'ils remplacent.}
Say too that it is known that $C(X)$ is of finite dimension over $\mathbb{Q}$,
(so that $A$ can also be expressed in terms of an equality of dimensions of
$C^{i}$ and $C^{n-1}_{\chi}$, which again proves it is independant of ), bu that
this is not known in car $p > 0$. Contrarily to what I hastily stated in my talk
(influenced from my recollections of the car $0$ case) it is not clear to me if
in car $p > 0$ the \struck{conditions} \struck{\ill{}} $A_{\chi}(X, \xi)$ and
$C_{\chi}(X, \xi)$ are independant of the polarization $\xi$ ; if you do not
find some proof of this independance, then the possible dependance should be
pointed out, as well as the fact that \struck{a priori} we do not have a proof
that $A$ to $D$ are independant of $\chi$. Of cour-

\page{44}

se, if the index theorem is proved for $X$, then $A_{\chi}(X, \xi) =
B_{\chi}(X)$ is again independant of the polarization, and analoguous remark for
$C_{\chi}(X, \xi)$.

When speaking about condition $C_{\chi}(X, \xi)$, emphasize at once its
stability properties by products (the proof I suggested works indeed)
specialization (with possible change of characteristics), hyperplane or more
generally linear sections. Give an extra proposition for the relations with the
property $A$, via a formal proposition as follows:

\textbf{Proposition} Condititions équivalentes sur $X$ (variété
\emph{polarisée}) :
\note{« Proposition » et « polarisée » sont soulignés à la machine ;
« Condititions » est de lui.}

\begin{enumerate}
\item[(i)] $C_{\chi}(X)$
\item[(ii)] $C_{\chi}(Y)$ et $A_{\chi}(X \times X)^{\equiv}$, \struck{\ill{}}
\item[(ii bis)] $C_{\chi}(Y)$ et $A_{\chi}(X \times X)^{0}$, où l'exposant
$^{0}$ signifie qu'on se borne à exprimer la condition $A$ pour
l'homomorphisme en dimension critique $H^{2n-2} \longrightarrow H^{2n+2}$.
\item[(iii)] $C_{\chi}(Y)$ et $A_{\chi}(X \times Y)$.
\item[(iii bis)] $C_{\chi}(Y)$ et $A_{\chi}(X \times Y)^{0}$, où l'exposant
$^{0}$ signifie qu'on se borne à exprimer la condition $A$ en dimension critique
$H^{(2n-1)-1} \longrightarrow H^{(2n-1)+1}$.
\item[(iv)] \struck{\ill{}} $C_{\chi}(Y)$, et pour tout $i \leq n-1$,
l'homomorphisme naturel $H^{i}(Y) \to H^{i}(X)$ inverse à gauche de
$\varphi^{i} : H^{i}(X) \to H^{i}(Y)$ (induit par $\Lambda_{X}\varphi_{*}$) est
induit par une classe de correspondance algébrique (induisant ce qu'elle veut
sur les autres $H^{j}(Y)$).
\item[(iv bis)] $C_{\chi}(Y)$, et pour $j \geq n+1$, l'homomorphisme naturel
$H^{j}(X) \to H^{j-2}(Y)$ inverse à droite de $\varphi_{j-2} : H^{j-2}(Y) \to
H^{j}(X)$ \struck{est} (induit par $\varphi^{*}\Lambda_{X}$) est induit par une
classe de correspondance algébrique (induisant ce qu'elle veut sur les autres
$H^{i}(X)$).
\end{enumerate}

\textbf{Corollaire} Ces conditions équivalent aussi à

\begin{enumerate}
\item[(v)] $A_{\chi}(X \times X)^{0} + A_{\chi}(Y \times Y)^{0} +
A_{\chi}(Z \times Z)^{0} + \quad$, où $X \supset Y \supset Z$ est une suite
décorissante de sections hyperplanes. \ill{}
\note{Deux mots ajoutés de sa main après « hyperplanes » ne se lisent pas. Les
indices $\chi$ du deuxième et du troisième terme ne sont pas encrés.}
\end{enumerate}

\page{45}

\begin{enumerate}
\item[(vi)] $A_{\chi}(X \times Y)^{0} + A_{\chi}(Y \times Z)^{0} + \dots$,
avec les mêmes notations.
\end{enumerate}

Of course, the products and hyperplane sections are endowed with the
polarizations stemming from the polarization on $X$. The conditions (v) and (vi)
have tthe slight interest that they allow to express the conjecture
$A(k) = C(k)$ in terms of $A(T)^{0}$ for every $T$ of even (resp.\ odd
dimension), where the upper $^{0}$ means that it is sufficient to look at what
happens in critical dimensions.

For thenproof of the proposition, I told you already the equivalence of (i) and
(ii), (ii bis). The equivalence of (iv) and (iv bis) is trivial by transposition,
they \struck{are of course implied by (i) because of the stabilities, because
$C_{\chi}(X)$ \ill{} $C_{\chi}(Y)$} imply (i) because $H^{2n-i}(X) \to H^{i}(X)$
is the compositum
\[
  H^{2n-i}(X) \longrightarrow H^{2n-i-2}(Y) \longrightarrow H^{i}(Y)
  \longrightarrow H^{i}(X)
\]
where the extreme arrows are the ones of (iv bis) and (iv) and the middle one is
induced by $\Lambda_{Y}^{(n-1)-i}$, and they are implied by (iii bis) because of
the formula
\[
  (\Lambda_{X}\varphi_{*})\,L_{Y} + L_{X}\,(\Lambda_{X}\varphi_{*})
  = (\varphi_{*}\Lambda_{Y}\varphi^{*} + \mathrm{id}_{X})\,\varphi_{*} .
\]
\note{L'exposant de $\Lambda_{Y}$ est frappé au-dessus de la ligne. Devant la
formule, deux caractères sont biffés.}
On the other hand (iii) $\Longrightarrow$ (iii bis) is trivial, and so is (i)
$\Longrightarrow$ (iii) because of the stabilities.

\marginal{N.B.\ $(\varphi^{*}\Lambda_{X})L_{X} + L_{Y}(\varphi^{*}\Lambda_{X})
= \varphi^{*}(\varphi_{*}\Lambda_{Y}\varphi^{*} + \mathrm{id}_{X})$}

For the list of the known facts, you can state that :

1) In arbitrary caracteristics, $C(X)$ is known if $\dim X \leq 2$, because more
generally, it is known that in arbitray dimension $n$, $H^{2n-1}(X) \to H^{1}(X)$
is induced by an algebraic correspondance class; also, in arbitray dimension, it
is known that $\pi_{0}$, $\pi_{2n}$, $\pi_{1}$, $\pi_{2n-1}$ are algebraic
(trivial for the first two, not quite trivial for the two next one's). If
$\dim X = 3$, it is not known however, even in car $0$, if $C(X)$ or only $D(X)$
hold, nor $A(X)$ and $B(X)$ in car.\ $p > 0$;

\marginal{also if for 1-cycles, $\tau$-équi.\ $=$ numér.\ équivalence …}

\page{46}

By the way, the fact that the $\pi_{i}$ for a surface are algebraic was pointed
out (Tate tells me) by Hodge in Algebraic correspondances between surfaces,
Proc.\ London Math.\ Soc.\ Seeies 2, Vol XLIV, 1938, p.226. It is rather striking
that this statement should not have struck the algebraic geometers more, and has
fallen into oblivion for nearly thirty years !
\note{« nearly » est ajouté de sa main au-dessus de la ligne.}

2) In car.$0$, $A(X)$ is known for $\dim X \leq 4$. But $A(X)^{0}$ is not known
if $\dim X = 5$; the first interesting case would be for a variety
$X \times Y$, $X$ of dim 3 and $Y$ a hyperplane section, as this would prove
$C(X)$, see above.

Thus the main problems arise already for 1-cycles on threefolds, and partially
even in car.$0$. Urged by Kleimann's question, I will look again at my old
scribbles on that subject (when I pretend to reduce the « strong » form of
Lefschetz to the « week » one). As for the suggestion I made on the phone, to
try to get any $X$ as birationally equivalent to a non singular $X'$, which is a
spcialisation of a non singular $X''$, itself birationally equivalent to a non
singular hypersurface --- this cannot work as Serre pointed out, because such an
$X$ would have to be simply connected ! Thus if one wants to reduce somehow to
the case of hypersurfaces, one will have to work also with singular ones, and
see how to reformulate for singular varieties the standard conjectures …

Sincerely your's

\section*{Lettre à Coates, 6 janvier 1966 (pages 47 à 49)}

\note{Second carbone, même présentation, la date « 6.1.1966 » de sa main en haut
à droite de la page 47. La première lettre porte « 4.1.67 » : on transcrit les
deux dates telles qu'elles sont écrites, sans les concilier.}

\page{47}

Dear Coates,

Here a few more comments to my talk on the conjectures. The following
proposition shows that the conjecture $C_{\chi}(X)$ is independant of the
choosen polarization, and has also some extra interest, in showing the part
playd by the fact that $H^{i}(X)$ should be « motive-theorietically »
isomorphic to its natural dual $H^{2n-i}(X)$ (as usual, I drop the twist for
simplicity).

\textbf{Proposition} The condition $C(X)$ is equivalent also to each of the
follwing conditions:
\note{L'indice de $C$ n'est pas encré ; partout ailleurs la lettre est
$\chi$. « each of » est frappé au-dessus de la ligne.}

\begin{enumerate}
\item[a)] $D_{\chi}(X)$ horlds, and for every $i < n$, there exists an
isomorphism $H^{2n-i}(X) \to H^{i}(X)$ which is algebraic (i.e.\ induced by an
algebraic correspondance class; we do not make any assertion on what it induces
in degrees different from $2n-i$).
\item[b)] For every endomorphism $H^{i}(X) \to H^{i}(X)$ which is algebraic, the
coefficients of the characteristic polynomial are rational, and for every
$i < n$, there exists an isomorphism $H^{2n-i}(X) \to H^{i}(X)$ which is
algebraic.
\end{enumerate}
\note{Dans a) et dans b), « $< n$ » est ajouté de sa main au-dessus du « $i$ ».}

Proof. I sketched already how $D_{\chi}(X)$ implies the fact that for an
algebraic endomorphism of $H^{i}(X)$, the coefficients of the characteristic
polynomial are rational, Therefore we know that a) implies b), and of course
$C_{\chi}(X)$ implies a). It remains to prove that b) implies $C_{\chi}(X)$. Let
$u : H^{2n-i}(X) \to H^{i}(X)$ be the given isomorphism which is algebraic, and
$v : H^{i}(X) \to H^{2n-i}(X)$ an algebraic isomorphism in the opposite
direction, induced by $L_{X}^{n-i}$. Then $uv = w$ is an automorphism of
$H^{i}(X)$ which is algebraic, and the Hamilton-Cayley formula
\[
  u^{b} - \sigma_{1}(w)\,u^{b-1} + \dots + \struck{(\pm 1)} (-1)^{b}\sigma_{b}(w) = 0
\]
(where the $\sigma_{i}(w)$ are the coefficients of the car.\ pol.\ of $w$) whos
that $w^{-1}$ is a linear combination of the $w^{i}$, with coefficients of the
type $\pm\,\sigma_{i}(w)/\sigma_{b}(w)$ (NB $b = \operatorname{rank} H^{i}$). The
assumption implies
\note{La formule de Hamilton-Cayley est écrite en $u$ alors que ses coefficients
sont ceux de $w$ ; on transcrit tel quel.}

\page{48}

that these coefficients are rational, which implies that $w^{-1}$ is algebraic,
and so is $w^{-1}u = v^{-1}$, which was to be proved.

NB In char.\ $0$, \struck{it is} the statement simplifies to: $C(X)$ equivalent
to the existence of algebraic isomorphisms $H^{2n-i}(X) \to H^{i}(X)$, (as the
preliminarxy condition in b) is then automatically satisfied). Maybe with some
extra care this can be proved too in arbitrary characteristics.

\textbf{Corollary} Assume $X$ and $X'$ satisfy condition $C_{\chi}$, and let
$u : H^{i}(X) \to H^{i+2j}(X')$ $(j \in \mathbb{Z})$ be an isomorphism which is
algebraic. Then $u^{-1}$ is algebraic.
\note{« $(j \in \mathbb{Z})$ » est ajouté de sa main au-dessus de la ligne.}

Indeed, the two spaces can be identified « algebraically » (both directions !)
to their dual, so taht the transpose of $u$ can be viewed as an isomorphism
$u' : H^{i+2j}(X') \to H^{i}(X)$. Thus $u'u$ is an algebraic automorphism $w$ of
$H^{i}(X)$, and by the previous argument we see that $w^{-1}$ is algebraic,
hence so is $u^{-1} = w^{-1}u'$.
\note{« algebraic » et le nom « $w$ » sont ajoutés de sa main au-dessus de la
ligne.}

As a consequence, we see that if $x \in H^{i}(X)$ is such that $u(x)$ is
algebraic ($i$ being now assumed to be even), then so is $x$. The same result
should hold in fact if $u$ is a monomorphism, the reason being that in this case
there should exist a left-inverse which is algebraic; this exists indeed in a
case like $H^{n-1}(X) \to H^{n-1}(Y)$ (where we take the left inverse
$\Lambda_{X}\varphi_{*}$). But to get it in general, it seems we need moreover
the Hodge index relation. (The complete yoga then being that we have the
category of motives which is semi-simple !). Without speeking of motives, and
staying down on earth, it would be nice to explain in the notes that $C(X)$
together with the index relation $I(X \times X)$ implies that the ring of
correspondance classes for $X$ is semi-simple, and how one deduces from this the
existence of left and right inverses as looked for above.

\page{49}

This could be given in an extra paragraph (which I did not really touch upon in
the talk), containing also the deduction of the Weil conjectures from the
conjectures $C$ and $A$.

A last and rather trivial remark is the following. Let's introduce
\struck{$x$} variants $A'_{\chi}(X)$ and $A''_{\chi}(X)$ as follows:

$A'_{\chi}(X)$ : if $2i \leq n-1$, any element $x$ of $H^{i}(X)$ whose image in
$H^{i}(Y)$ is algebraic, is algebraic.
\note{« $x$ » est ajouté de sa main au-dessus de « element ».}

$A''_{\chi}(X)$ : if $2i \geq n-1$, any element \struck{\ill{}} of $H^{i+2}(X)$
is \struck{\ill{}} the image of an algebraic element of $H^{i}(Y)$.
\note{La ligne est refaite : « algebraic » est ajouté de sa main au-dessus, et
tout ce qui sépare « element » de « $H^{i+2}(X)$ », puis le début de la ligne
suivante, sont biffés à la machine.}

Let us consider also the specifications $A'_{\chi}(X)^{0}$ and
$A''_{\chi}(X)^{0}$, where we restrict to the crittval dimensions $2i = n-1$ if
$n$ odd, $2i = n-2$ if $n$ even. All these conditions are in the nature of
« week » Lefschetz relations, and they are trivially implied by $A_{\chi}(X)$
resp.\ $C_{\chi}(Y)$ (in the first case, applying $\varphi$ we see that
$L_{X}x$ is algebraic; in the second, we take $y = \Lambda_{Y}\varphi^{+}(x)$).
The remark then is that these pretendedly « week » variants in fact imply the
full Lefschetz relations for algebraic cycles, namely:

\textbf{Proposition} $C_{\chi}(X)$ is equivalent to the conjunction
$C_{\chi}(Y) + A'_{\chi}(X \times X)^{0} + A''_{\chi}(X \times X)^{0}$, hence
(by induction) also to the conjunction of the conditions $A'^{\,0}_{\chi}$ and
$A''^{\,0}_{\chi}$ for all of the varieties $X \times X$, $Y \times Y$,
$Z \times Z$ … Analoguous statement with $X \times Y$, $Y \times Z$ etc
instead of $X \times X$, $Y \times Y$ etc.

This comes from the remark that \struck{\ill{}} $A_{\chi}(X)^{0}$ follows from
the conjuction of $A'_{\chi}(X)^{0}$ and $A''_{\chi}(X)^{0}$, as one sees by
decomposing $L_{X}^{2} : H^{2m-2}(X) \to H^{2m+2}(X)$ into
\[
  H^{2m-2}(X) \xrightarrow{\ \varphi^{*}\ } H^{2m-2}(Y)
  \xrightarrow{\ \varphi_{*}\ } H^{2m}(X)
  \xrightarrow{\ L_{X}\ } H^{2m+2}(X)
\]
if $\dim X = 2m$ is even, and $H^{2m+1-1}(X) \to H^{2m+1+1}(X)$ into
\[
  H^{2m}(X) \xrightarrow{\ \varphi^{*}\ } H^{2m}(Y)
  \xrightarrow{\ \varphi_{*}\ } H^{2m+2}(X)
\]
if $\dim X = 2m+1$ is odd.
\note{Les noms des flèches, $\varphi^{*}$, $\varphi_{*}$ et $L_{X}$, sont écrits
de sa main au-dessus de chacune. Une phrase entière est biffée à la machine
après « the remark that ».}

Sincerely your's

\section*{Lemme : un anneau d'endomorphismes qui contient $L$ (pages 50, 52 et 54)}

\note{Trois feuillets de sa main, écrits au dos de la « Tribu » n° 65 de
Bourbaki. C'est vraisemblablement la suite des « old scribbles » que la lettre de
la page 46 promet de reprendre : la réduction de la forme forte de Lefschetz à
la forme faible. $X$ est de dimension $n$ et $H^{*}(X)$ sa cohomologie totale ;
$\zeta$ est l'opérateur donné, que la conclusion identifie à $L$, et $p_{i}$,
$P^{i}$ sont les projections sur $H^{i}$ et sur sa partie primitive.}

\page{50}

\textbf{Lemme} Soit $A$ un anneau d'endomorphismes du gpe abélien $H^{*}(X)$,
stable par composantes homogènes (\uncertain{deg.} \emph{tot.}), par
transposition. Supposons $\zeta \in A$, et \struck{\ill{}} que pour tout
$0 \leq i < n$, il existe $u_{i} \in A$, tel que
\[
  u_{i}\,\zeta^{n-i} : H^{i}(X) \longrightarrow H^{2n-i}(X) \longrightarrow H^{i}(X)
\]
soit l'identité. Alors $A$ contient $L$, \struck{\ill{}} et les projections sur
les facteurs élémentaires $\zeta^{j}P^{i}$ ($i \leq \uncertain{n}$, et
\struck{\ill{}} $j \leq n-i$).
\note{« Supposons $\zeta \in A$, et » est ajouté en interligne à droite ; une
seconde insertion, au-dessus de « tel que », ne se lit qu'en partie :
\ill{} $2n - 2(n-i)$). « tel que » est encadré.}

\marginal{(\ill{})}

\textbf{Corollaire} Conditions équivalentes :

\begin{enumerate}
\item[(i)] $L$ algébrique
\item[(ii)] \ill{} et $1 \leq i \leq n$, il existe une corr.\ \emph{alg.}
induisant $H^{2n-i}(X) \to H^{i}(X)$ inverse de $\uncertain{\zeta^{i}}$.
\end{enumerate}

\struck{(iii)} S'il en est ainsi, alors \ill{} les projections sur les facteurs
élémentaires de $H^{*}(X)$, sont algébriques, et les projections
$H^{*}(X) \to H^{i}(X)$, aussi.

\emph{Dém.\ du lemme.} Comme $A$ stable par transposition, il suffit de regarder
les $\zeta^{j}P^{i}$, pour $i + 2j \leq n$. Soit $k = i + 2j$, --- procédons par
récurrence sur $k$, trivial si $k < 0$.

\page{52}

Supposons $k \geq 0$, et le \ill{} \ill{} pour $k' < k$. On sait déjà que les
$p_{i} : H^{\uncertain{*}}(X) \to H^{i}(X)$ sont $\in \mathrm{Alg}$ pour
$i < k$, donc (transp.) pour $i > 2n-k$. Soit
\[
  T_{k} = \mathrm{id} - \sum_{i<k} p_{i} - \sum_{j>2n-k} p_{j} ;
\]
donc $T_{k} \in A$, plus \uncertain{gén.}\ $T_{k'} \in A$ si $k' \leq k$. Alors
$\zeta^{n-k}T_{k}$ est nul sur $H^{j}$, $j \neq k$, et induit
$H^{k} \xrightarrow{\ \zeta^{n-k}\ } H^{2n-k}$, donc
$u_{k}\,\zeta^{n-k}\,T_{k} = p_{k}$.

D'autre part, si $i + 2j = k$, $j > 0$, les projections sur $\zeta^{j}p^{i}$
sont aussi $\in A$, comme \ill{} voit \struck{aussitôt} par récurrence
ascendante sur $i$ (ou descendante sur $j$) ainsi : si $i < 0$ rien à prouver.
Si $i$ \ill{} \struck{\ill{}} pour les $i' < i$, \ill{} $\zeta^{n-k+j}T_{k}$,
\struck{\ill{}} qui \ill{} nul sur $H^{\alpha}$ pour $\alpha \neq k$, et induit
$H^{k} \xrightarrow{\ \zeta^{n-k+j}\ } H^{2n-k+2j}$, \ill{} est nul sur les
facteurs $\zeta^{j'}p^{i'}$ pour $j' < j$, (et sur les facteurs
$\zeta^{j'}p^{i'}$ avec $j' \geq j$ (i.e.\ $i' \leq i$), \struck{coïncide avec
$1$}) \ill{} $\zeta^{j}H^{k-2j}$, \ill{} $\zeta^{n-(k-2j)}$. Donc l'homom.
\[
  H^{k-2j} \xrightarrow{\ \zeta^{j}\ } \zeta^{j}H^{k-2j}
\]
et $u_{k-2j}\,\zeta^{n-k+j}\,T_{k}$ \ill{} ; nul sur $N$, et sur
$\zeta^{j}H^{k-2j}$ \ill{} induisant $H^{k-2j} \xrightarrow{\ \sim\ }
\zeta^{j}H^{k-2j}$. Donc \ill{}
\note{Cette page est écrite très vite : les formules portent, la prose qui les
relie souvent non. Une insertion en interligne, au-dessus de la flèche
$H^{k} \to H^{2n-k+2j}$, se termine par « des mêmes $N$, » et ne se lit pas
autrement ; $N$ n'est défini nulle part sur ces trois feuillets.}

\page{54}

$\zeta^{j}u_{k-2j}\,\zeta^{n-k+j}\,T_{k}$ n'est \uncertain{autre} que la
projection de $H^{\alpha}$ sur $\zeta^{j}H^{k-2j}$, \ill{} des
$\zeta^{j'}p^{i'}$ pour $i' + 2j' = k$, $j' \geq j$ ; \struck{car} comme \ill{}
les projections \struck{sur} partielles sur $\zeta^{j'}p^{i'}$, $j' > j$, sont
$\in A$, \ill{} \ill{} \ill{} $=$ \ill{} la projection sur $\zeta^{j}p^{i}$, par
différence.

Ceci prouve le lemme.

\end{document}
